\chapter{Infinite \texorpdfstring{\(q\)}{q}-Pochhammer inverse geometry}
\label{chap:infinite-q-pochhammer-inversion}

\section{The infinite product: analytic domain and certified tails}
\label{sec:ip-normal-convergence}
\label{sec:infinite-normal-convergence}

For \(|q|<1\), put
\[
 P(a,q):=(a;q)_\infty=\prod_{j=0}^{\infty}(1-aq^j).
\]
The domain condition is part of the definition.  Values on \(|q|=1\), radial
limits toward that boundary, and exterior reciprocal normalizations are
different objects and will be treated separately.

\begin{theorem}[Normal convergence, zeros, and derivative kernels]
\label{thm:ip-normal-convergence}
Let \(0<\rho<1\) and \(R>0\).  Uniformly for \(|q|\leq\rho\), \(|a|\leq R\),
and \(N\geq0\),
\begin{equation}
 \left|\prod_{j=N}^{\infty}(1-aq^j)-1\right|
 \leq \exp\left(\frac{R\rho^N}{1-\rho}\right)-1.
 \label{eq:ip-product-tail}
\end{equation}
If \(R\rho^N<1\), the logarithm obtained by continuation from \(a=0\)
satisfies
\begin{equation}
 \left|\sum_{j=N}^{\infty}\log(1-aq^j)\right|
 \leq
 \frac{R\rho^N}{(1-\rho)(1-R\rho^N)}.
 \label{eq:ip-log-tail}
\end{equation}
Consequently \(P\) is holomorphic on
\(\{(a,q)\in\ComplexNumbers^2:|q|<1\}\), and is entire in \(a\) for fixed \(q\).
For \(0<|q|<1\), its zeros as a function of \(a\) are exactly
\(a=q^{-m}\), \(m\geq0\), and they are simple.  Away from those zeros,
\begin{align}
 \partial_a\log P(a,q)
  &=-\sum_{j\geq0}\frac{q^j}{1-aq^j},
  \label{eq:ip-a-logder}\\
 \partial_a^2\log P(a,q)
  &=-\sum_{j\geq0}\frac{q^{2j}}{(1-aq^j)^2},
  \label{eq:ip-a-logsecond}\\
 \partial_q\log P(a,q)
  &=-a\sum_{j\geq1}\frac{j q^{j-1}}{1-aq^j}.
  \label{eq:ip-q-logder}
\end{align}
All three series converge locally uniformly on the indicated zero-free set.
\end{theorem}

\begin{proof}
Put \(u_j=-aq^j\).  Expanding finite products and applying the triangle
inequality gives
\[
 \left|\prod_{j=N}^{M}(1+u_j)-1\right|
 \leq\prod_{j=N}^{M}(1+|u_j|)-1
 \leq\exp\left(\sum_{j=N}^{M}|u_j|\right)-1.
\]
Since \(\sum_{j\geq N}|u_j|\leq R\rho^N/(1-\rho)\), the finite tails form a
uniformly Cauchy family and their limit obeys \cref{eq:ip-product-tail}.  If
\(R\rho^N<1\), then
\[
 |\log(1-z)|\leq\sum_{k\geq1}\frac{|z|^k}{k}
 \leq\frac{|z|}{1-|z|}
\]
for \(|z|<1\).  With \(|aq^j|\leq R\rho^N\rho^{j-N}\), summation proves
\cref{eq:ip-log-tail}.

Locally uniform limits of the finite holomorphic products are holomorphic.
If no factor vanishes, choose \(N\) so that \(|aq^j|<1/2\) for \(j\geq N\).
The tail is the exponential of an absolutely convergent logarithmic sum and
is therefore nonzero; hence the only zeros are the factor zeros.  At
\(a=q^{-m}\), exactly one factor vanishes, so the zero is simple.  On a compact
zero-free set, the denominators are separated from zero for the finite head,
while the tails in \cref{eq:ip-a-logder,eq:ip-a-logsecond,eq:ip-q-logder} are
bounded by convergent geometric series with at most a polynomial factor in
\(j\).  Normal convergence therefore permits termwise differentiation and
proves all three identities.
\end{proof}

\paragraph{Exact fixed-nome Lean subclaims and their boundary.}
The generic formal object is \lean{Fabius.qPochhammerInfIn}, the total
\texttt{tprod} in an arbitrary topological commutative ring, with defining
theorem \lean{Fabius.qPochhammerInfIn_eq_tprod}.  In a normed ring with
\(\lVert1\rVert=1\), the elementary lemmas are
\lean{Fabius.summable_norm_mul_pow},
\lean{Fabius.one_sub_ne_zero_of_norm_lt_one}, and
\lean{Fabius.norm_mul_pow_self_lt_one}; the finite self-product result
\lean{Fabius.finiteQPochhammerIn_self_ne_zero} additionally uses a
commutative ring with no zero divisors.

For \(\lVert q\rVert<1\) in a complete normed commutative ring with
\(\lVert1\rVert=1\), the seven exact convergence and factorization theorems
are
\lean{Fabius.multipliable_one_sub_mul_pow_of_norm_lt_one},
\lean{Fabius.hasProd_qPochhammerInfIn},
\lean{Fabius.tendsto_finiteQPochhammerIn_qPochhammerInfIn},
\lean{Fabius.qPochhammerInfIn_eq_finite_mul_shift},
\lean{Fabius.qPochhammerInfIn_succ_shift},
\lean{Fabius.qPochhammerInfIn_eq_factor_mul}, and
\lean{Fabius.qPochhammerInfIn_dissection}; only the last also assumes
\(0<r\).  With a multiplicative norm, the exact zero block is
\lean{Fabius.qPochhammerInfIn_ne_zero},
\lean{Fabius.qPochhammerInfIn_eq_zero_iff}, and
\lean{Fabius.qPochhammerInfIn_self_ne_zero}.  The middle theorem is the raw
factor equation \(aq^j=1\), so it includes \(q=0\).

The unconditional normed-field support theorems are
\lean{Fabius.qPochhammerInfIn_eq_tprod_smul} and
\lean{Fabius.isBigO_one_sub_sub_one};
\lean{Fabius.summable_norm_pow_of_norm_lt_one} additionally assumes
\(\lVert q\rVert<1\), while
\lean{Fabius.differentiable_finiteQPochhammerIn} needs only a nontrivially
normed field.  In a complete normed field, \(\lVert q\rVert<1\) and
\(q\ne0\) give the reciprocal
zero enumeration
\lean{Fabius.qPochhammerInfIn_eq_zero_iff_exists_inv_pow}; the same strict
contraction together with local compactness is assumed by
\lean{Fabius.hasProdLocallyUniformly_qPochhammerInfIn} and
\lean{Fabius.continuous_qPochhammerInfIn}.  The finite cofactor identity
\lean{Fabius.pow_sq_mul_finiteQPochhammerIn_inv_pow_self} is algebraic over a
field with \(q\ne0\).

Over \(\ComplexNumbers\), the final four theorems are
\lean{Fabius.differentiable_qPochhammerInfIn},
\lean{Fabius.hasDerivAt_qPochhammerInfIn_of_mul_pow_eq_one},
\lean{Fabius.hasDerivAt_qPochhammerInfIn_inv_pow}, and
\lean{Fabius.deriv_qPochhammerInfIn_inv_pow_ne_zero}.  The raw derivative
theorem assumes \(a_0q^j=1\) and permits \(q=0,j=0\); the reciprocal-power
formula and its nonvanishing theorem assume \(q\ne0\) and give exactly
\cref{eq:ip-zero-derivative}.  The final theorem literally proves the
displayed coefficient nonzero, while the preceding \texttt{HasDerivAt}
theorem identifies that coefficient with the derivative.

The finite parent module \path{QPochhammerDissection.lean} has no public
definition and exactly the two commutative-ring theorems
\lean{Fabius.finiteQPochhammerIn_dissection} and
\lean{Fabius.finiteQPochhammerIn_dissection_remainder}.  They require no
condition on \(q\); the first is total in natural \(r,n\), and the second
assumes exactly \(u\le r\), including \(u=r\).  The positive-modulus
infinite identity listed above is the limit-level companion.

The infinite declarations listed above exhaust the one definition and
twenty-seven theorems of
\path{QPochhammerInfinite.lean}.  They formalize the fixed-\(q\) convergence,
entireness, zero, finite-shift, residue-dissection, and zero-derivative
subclaims used in this chapter.  They do \emph{not} formalize the
uniform-in-\(q\) numerical tail bounds
\cref{eq:ip-product-tail,eq:ip-log-tail}, joint holomorphy in \((a,q)\), or the
three logarithmic-derivative series
\cref{eq:ip-a-logder,eq:ip-a-logsecond,eq:ip-q-logder}; consequently the
compound \Cref{thm:ip-normal-convergence} remains a human-proved frontier
result.  Nor does this generic module state analytic order.  That separate
conclusion for \lean{Fabius.complexQPochhammerInf}, including \(q=0\), remains
\lean{Fabius.analyticOrderAt_complexQPochhammerInf_of_eq_zero} in
\path{QPochhammerEntire.lean}.

\begin{lemma}[Lambert logarithm]
\label{lem:ip-lambert-log}
For \(|a|<1\) and \(|q|<1\),
\begin{equation}
 \log(a;q)_\infty
 =-\sum_{m\geq1}\frac{a^m}{m(1-q^m)},
 \label{eq:ip-lambert-log}
\end{equation}
where the logarithm is the branch equal to zero at \(a=0\).  The double
series used in its derivation is absolutely and locally uniformly convergent.
\end{lemma}

\begin{proof}
For each factor,
\(\log(1-aq^j)=-\sum_{m\geq1}a^mq^{jm}/m\).  On compact subdiscs
\(|a|\leq\alpha<1\), \(|q|\leq\rho<1\),
\[
 \sum_{j\geq0}\sum_{m\geq1}\frac{|a|^m|q|^{jm}}m
 \leq\sum_{m\geq1}\frac{\alpha^m}{m(1-\rho^m)}
 \leq\frac{-\log(1-\alpha)}{1-\rho}<\infty.
\]
Thus Fubini's theorem applies, and summing
\(\sum_{j\geq0}q^{jm}=(1-q^m)^{-1}\) gives
\cref{eq:ip-lambert-log}.
\end{proof}

\section{Argument inversion and the complete real sheet atlas}
\label{sec:ip-argument-atlas}

\begin{theorem}[Principal inverse, all gaps, and exact zero normalization]
\label{thm:ip-real-atlas}
Fix \(0<q<1\), write \(P_q(a)=(a;q)_\infty\), and put
\(r_m=q^{-m}\).
\begin{enumerate}
\item \(P_q\) is positive and strictly decreasing from \(+\infty\) to \(0\)
  on \((-\infty,1)\).  It therefore has a unique real-analytic inverse
  \(\mathcal A_q:(0,\infty)\to(-\infty,1)\).
\item Every \(r_m\) is a simple zero, with
  \begin{equation}
   d_m:=P_q'(r_m)=(-1)^{m+1}q^{-m(m-1)/2}
   (q;q)_m(q;q)_\infty.
   \label{eq:ip-zero-derivative}
  \end{equation}
  Hence its inverse germ is \(a_m(y)=r_m+y/d_m+\BigO(y^2)\).
\item On the gap \(I_m=(r_m,r_{m+1})\), the sign of \(P_q\) is
  \(\sigma_m=(-1)^{m+1}\).  There is exactly one critical point
  \(c_m\in I_m\), and \(P_q''(c_m)\ne0\).  If
  \(M_m=|P_q(c_m)|\), then \(I_m\) contains two roots of \(P_q(a)=y\)
  when \(0<\sigma_my<M_m\), one double root when
  \(\sigma_my=M_m\), and no root when
  \(\sigma_my<0\) or \(\sigma_my>M_m\).  At \(y=0\), its two endpoint
  roots are \(r_m,r_{m+1}\).
\end{enumerate}
The principal inverse is strictly decreasing and convex, and, with
\[
 S_k(a)=\sum_{j\geq0}\frac{q^{kj}}{(1-aq^j)^k},
\]
it satisfies
\begin{align}
 \mathcal A_q'(y)&=-\frac1{yS_1(\mathcal A_q(y))},
 \label{eq:ip-principal-first}\\
 \mathcal A_q''(y)&=
 \frac{S_1(\mathcal A_q(y))^2-S_2(\mathcal A_q(y))}
 {y^2S_1(\mathcal A_q(y))^3}>0.
 \label{eq:ip-principal-second}
\end{align}
\end{theorem}

\begin{proof}
For \(a<1\), all factors are positive and
\(\partial_a\log P_q<0\) by \cref{eq:ip-a-logder}.  As \(a\uparrow1\),
\[
 P_q(a)=(1-a)(aq;q)_\infty
 \sim(q;q)_\infty(1-a)\longrightarrow0,
\]
where normal convergence gives the limit of the second factor.  If \(a=-x\),
then \(P_q(-x)=\prod_{j\geq0}(1+xq^j)\geq1+x\to\infty\).
This proves the first bijection and, since its derivative never vanishes,
real analyticity of its inverse.

To compute the derivative at \(r_m\), remove the vanishing factor:
\[
 P_q'(r_m)=-q^m\prod_{j\ne m}(1-q^{j-m}).
\]
The indices \(j<m\) contribute
\((-1)^mq^{-m(m+1)/2}(q;q)_m\); the indices \(j>m\) contribute
\((q;q)_\infty\).  This proves \cref{eq:ip-zero-derivative} and the inverse
germ.

On every compact subset avoiding the zeros, \cref{eq:ip-a-logder} may be
written
\[
 \frac{P_q'(a)}{P_q(a)}=\sum_{j\geq0}\frac1{a-r_j},\qquad
 \left(\frac{P_q'}{P_q}\right)'(a)
 =-\sum_{j\geq0}\frac1{(a-r_j)^2}<0.
\]
In \(I_m\), the first sum tends to \(+\infty\) at the left endpoint and to
\(-\infty\) at the right endpoint.  It therefore has one and only one zero.
At that point \(P_q''/P_q\) equals the strictly negative second sum, so the
critical point is nondegenerate.  There are \(m+1\) negative factors in the
gap, giving the sign \(\sigma_m\).  The logarithmic derivative has opposite
signs on the two sides of \(c_m\); hence the absolute value rises strictly
from zero to \(M_m\) and falls strictly to zero.  The root count follows.

Implicit differentiation gives \cref{eq:ip-principal-first}.  Since
\[
 S_1(a)^2-S_2(a)
 =2\sum_{0\leq i<j}
 \frac{q^{i+j}}{(1-aq^i)(1-aq^j)}>0
\]
for \(a<1\), a second differentiation gives
\cref{eq:ip-principal-second} and proves convexity.
\end{proof}

\paragraph{Lean boundary.}
The zero and derivative subclaim is exact through
\lean{hasDerivAt_qPochhammerInfIn_inv_pow} and
\lean{deriv_qPochhammerInfIn_inv_pow_ne_zero}.  The real inverse germs,
monotonicity and convexity, uniqueness of the gap critical point, and sheet
counts in \cref{thm:ip-real-atlas} remain human-proved.

\subsection{Endpoint and central inverse jets}

Euler's expansion
\[
 (a;q)_\infty=\sum_{m\geq0}
 \frac{(-1)^mq^{\binom m2}}{(q;q)_m}a^m
\]
and Lagrange inversion give, with \(s=1-y\),
\begin{align}
 \mathcal A_q(y)={}&(1-q)s+\frac{q(1-q)}{1+q}s^2 \notag\\
 &+\frac{q^2(1-q)(q^2+q+2)}
 {(1+q)^2(q^2+q+1)}s^3+\BigO(s^4),
 \label{eq:ip-unit-target-jet}\\
 [s^m]\mathcal A_q(1-s)
 ={}&\frac1m[z^{m-1}]
 \left(\frac{z}{1-(z;q)_\infty}\right)^m.
 \label{eq:ip-unit-target-all-order}
\end{align}
This follows exactly as in \cref{fp-unit-inverse}: the forward linear
coefficient is \(1/(1-q)\), and direct coefficient comparison gives the three
displayed terms.

At the other endpoint, let
\[
 C_q=(q;q)_\infty,\qquad
 L_1(q)=\sum_{j\geq1}\frac{q^j}{1-q^j},\qquad v=\frac y{C_q}.
\]
Writing \(a=1-\delta\) gives
\[
 (a;q)_\infty=C_q\delta\bigl(1+L_1(q)\delta+\BigO(\delta^2)\bigr).
\]
Reversion therefore yields
\begin{equation}
 \mathcal A_q(y)=1-v+L_1(q)v^2+\BigO(v^3).
 \label{eq:ip-zero-target-jet}
\end{equation}
The radius of either inverse series is bounded by the nearest singular value
encountered by its chosen sheet; such a singular value may be critical or
asymptotic, and need not be the numerically nearest critical value belonging
to some other sheet.

\subsection{Remote real critical points}

\begin{theorem}[Remote-gap asymptotics]
\label{thm:ip-remote-gaps}
Let \(m=k+1\).  As \(k\to\infty\),
\begin{align}
 c_k(q)&=q^{-m}\left(1-\frac1{m+1}+\BigO(m^{-2})\right),
 \label{eq:ip-remote-location}\\
 |P_q(c_k)|
 &= (q;q)_\infty^2q^{-m(m+1)/2}
 \frac{m^m}{(m+1)^{m+1}}\left(1+\BigO(m^{-1})\right)
 \label{eq:ip-remote-value}\\
 &=\frac{(q;q)_\infty^2}{\EulerE(m+1)}
 q^{-m(m+1)/2}\left(1+\BigO(m^{-1})\right).
 \notag
\end{align}
The error constants are locally uniform for \(q\) in compact subsets of
\((0,1)\).
\end{theorem}

\begin{proof}
Parameterize the gap by \(a=q^{-m}(1-\varepsilon)\),
\(0<\varepsilon<1-q\).  Separating the \(m\) factors before the right endpoint
gives the exact identity
\begin{align}
 P_q(a)
 &=(-1)^mq^{-m(m+1)/2}
 \varepsilon(1-\varepsilon)^mG_m(\varepsilon),
 \label{eq:ip-remote-factorization}\\
 G_m(\varepsilon)
 &=\left(\frac q{1-\varepsilon};q\right)_m
 ((1-\varepsilon)q;q)_\infty.
 \notag
\end{align}
If \(\varepsilon=\BigO(m^{-1})\), logarithmic differentiation of the two
products shows uniformly that
\[
 \partial_\varepsilon\log G_m(\varepsilon)=\BigO(1).
\]
Indeed, the resulting summands are bounded by constant multiples of
\(q^j/(1-q^j)\), whose sum converges.  The same logarithmic estimate, together
with the omitted tail \(\sum_{j>m}q^j=\BigO(q^m)\), gives
\[
 G_m(\varepsilon)=(q;q)_\infty^2(1+\BigO(m^{-1})).
\]
It remains to verify that the unknown extremum lies in this range.  The
logarithmic derivative of \cref{eq:ip-remote-factorization} is positive as
\(\varepsilon\downarrow0\).  At
\(\varepsilon=2/(m+1)\), which lies inside the gap for all sufficiently large
\(m\), the first two terms below have size \(-m/2+\BigO(1)\), while the
already established \(G_m'/G_m=\BigO(1)\) estimate applies there.  The
derivative is therefore negative.  Uniqueness of the gap extremum gives
\(\varepsilon_*<2/(m+1)\), so \(\varepsilon_*=\BigO(m^{-1})\) without
circular reasoning.

At that unique extremum, logarithmic differentiation of
\cref{eq:ip-remote-factorization} now gives
\[
 0=\frac1\varepsilon-\frac m{1-\varepsilon}
   +\partial_\varepsilon\log G_m(\varepsilon).
\]
The bounded last term perturbs the zero of the first two terms by
\(\BigO(m^{-2})\), so
\(\varepsilon=(m+1)^{-1}+\BigO(m^{-2})\).  Substitution proves
\cref{eq:ip-remote-location,eq:ip-remote-value}; finally
\((1-1/(m+1))^m=\EulerE^{-1}(1+\BigO(m^{-1}))\).
\end{proof}

\section{Large targets and the log-periodic principal inverse}
\label{sec:ip-log-periodic}

\begin{theorem}[Theta-modulated growth and its inverse]
\label{thm:ip-log-periodic}
Let \(q=\EulerE^{-\tau}\), \(\tau>0\).  For \(x=\EulerE^L\geq1\), write uniquely
\(x=q^{-m}r\), \(m\in\NaturalNumbers\), \(1\leq r<q^{-1}\), and
\(\rho=\log r\in[0,\tau)\).  Define on that interval
\begin{equation}
 \Psi_q(\rho)=
 \log\left((-\,\EulerE^\rho;q)_\infty(-q\EulerE^{-\rho};q)_\infty\right)
 -\frac\rho2-\frac{\rho^2}{2\tau},
 \label{eq:ip-phase}
\end{equation}
and extend it \(\tau\)-periodically.  Then \(\Psi_q\) is smooth and
\begin{equation}
 \log(-\EulerE^L;q)_\infty
 =\frac{L^2}{2\tau}+\frac L2+\Psi_q(L\bmod\tau)
 +\BigO(\EulerE^{-L}).
 \label{eq:ip-log-periodic-forward}
\end{equation}
If \(Y=(-x;q)_\infty\to\infty\), set
\[
 R=\sqrt{2\tau\log Y},\qquad
 \rho_0=(R-\tau/2)\bmod\tau.
\]
Then
\begin{equation}
 \log x=R-\frac\tau2+
 \frac{\tau\bigl(\tau/8-\Psi_q(\rho_0)\bigr)}R
 +\BigO((\log Y)^{-1}).
 \label{eq:ip-log-periodic-inverse}
\end{equation}
\end{theorem}

\begin{proof}
The finite shift is exact:
\[
 (-q^{-m}r;q)_\infty
 =r^mq^{-m(m+1)/2}(-q/r;q)_m(-r;q)_\infty.
\]
Uniformly for \(1\leq r\leq q^{-1}\), replacing the finite product by
\((-q/r;q)_\infty\) has relative error \(1+\BigO(q^m)\), by
\cref{eq:ip-product-tail}.  Since \(L=m\tau+\rho\),
\[
 m\rho+\frac{\tau m(m+1)}2
 =\frac{L^2}{2\tau}+\frac L2-\frac{\rho^2}{2\tau}-\frac\rho2.
\]
Taking logarithms proves \cref{eq:ip-log-periodic-forward}.  Normal
convergence permits all \(\rho\)-derivatives on the compact fundamental
interval.  Under \(\rho\mapsto\rho+\tau\), the product inside the logarithm
is multiplied by \(\EulerE^{\rho+\tau}\), exactly the change of the two subtracted
terms, so the extension is smooth and periodic.

Completing the square rewrites the forward equation as
\[
 2\tau\log Y=(L+\tau/2)^2+
 2\tau\bigl(\Psi_q(\rho)-\tau/8\bigr)+\BigO(\EulerE^{-L}).
\]
Taking the positive square root first gives \(L=R-\tau/2+\BigO(R^{-1})\).
Insert this once into the bounded correction.  The exact phase differs from
\(\rho_0\), on the periodic circle, by \(\BigO(R^{-1})\); smoothness of
\(\Psi_q\) changes the \(R^{-1}\) correction by only \(\BigO(R^{-2})\).
Since \(R^{-2}=\BigO((\log Y)^{-1})\), this is
\cref{eq:ip-log-periodic-inverse}.
\end{proof}

\section{\texorpdfstring{Inversion in the base at \(q=0\) and on the real interval}{Inversion in the base at q=0 and on the real interval}}
\label{sec:ip-real-base}

\begin{theorem}[Real base atlas and the stable small-base jet]
\label{thm:ip-real-base}
Fix real \(a\).
\begin{enumerate}
\item If \(0<a<1\), \(q\mapsto(a;q)_\infty\) decreases strictly on
  \([0,1)\) from \(1-a\) to \(0\).
\item If \(a<0\), it increases strictly from \(1-a\) to \(+\infty\).
\item If \(a=0\) it is identically \(1\), and if \(a=1\) it is identically
  \(0\).
\item If \(a>1\), it has the distinct simple base zeros
  \(q_j=a^{-1/j}\), \(j\geq1\), accumulating at \(1\); hence only
  interval-by-interval base inverses are possible.
\end{enumerate}
For \(a\notin\{0,1\}\), set
\(\eta=((1-a)-y)/(a(1-a))\).  The branch through \(q=0\) is
\begin{align}
 q={}&\eta-\eta^2+(1+a)\eta^3-(1+4a)\eta^4 \notag\\
 &+(1+11a+3a^2)\eta^5-(1+26a+22a^2)\eta^6
 +\BigO(\eta^7).
 \label{eq:ip-small-base-jet}
\end{align}
Every coefficient is obtained by finite exact arithmetic: through order
\(\eta^M\), it agrees with every finite product of order \(n>M\).
\end{theorem}

\begin{proof}
On \(0<q<1\), \cref{eq:ip-q-logder} has sign \(-\Sign(a)\) whenever \(a<1\)
and \(a\ne0\).  If \(0<a<1\), then for any \(M\), and \(q\) sufficiently
close to one, the first \(M\) nonconstant factors are at most
\(1-a/2<1\); the remaining factors lie in \((0,1)\).  Thus the product is at
most \((1-a/2)^M\), and its limit is zero.  If \(a<0\), the first \(M\)
factors are at least \(1+|a|/2\), while all remaining factors exceed one, so
the limit is \(+\infty\).  The values at \(q=0\) and the two constant cases
are immediate.  For \(a>1\), exactly the \(j\)-th factor vanishes at
\(q_j\), proving simplicity and the asserted accumulation.

Near \(q=0\), normal convergence makes the product analytic and
\(\partial_qP(a,0)=-a(1-a)\ne0\), so the analytic inverse exists.  The
coefficient of \(q^M\) uses only factors with index at most \(M\); hence it
equals the stable finite-product coefficient.  Direct subset expansion gives
\[
 \frac{1-(aq;q)_\infty}{a}
 =q+q^2+(1-a)q^3+(1-a)q^4+(1-2a)q^5
 +(1-2a+a^2)q^6+\BigO(q^7).
\]
Substitution of \(q=\eta+\sum_{j=2}^{6}c_j\eta^j\), followed by comparison of
successive powers, gives
\[
 (c_2,c_3,c_4,c_5,c_6)
 =(-1,\,1+a,\,-1-4a,\,1+11a+3a^2,\,-1-26a-22a^2),
\]
which is \cref{eq:ip-small-base-jet}.  Because the subset calculation through
degree \(M\) uses no factor beyond index \(M\), it also proves the stated
finite exactness mechanism.
\end{proof}

\section{\texorpdfstring{Fixed-argument inversion as \(q\to1\)}{Fixed-argument inversion as q approaches 1}}
\label{sec:ip-fixed-a-qone}

\begin{theorem}[Uniform all-order fixed-argument expansion]
\label{thm:ip-fixed-a-expansion}
Let \(q=\EulerE^{-t}\), where \(t\to0\) in a closed subsector of
\(\Re t>0\).  Uniformly for \(|a|\leq\rho<1\),
\begin{equation}
 \log(a;\EulerE^{-t})_\infty
 \sim-\frac{\Polylogarithm_2(a)}t+\frac12\log(1-a)
 -\sum_{r\geq1}\frac{B_{2r}}{(2r)!}
 \Polylogarithm_{2-2r}(a)t^{2r-1}.
 \label{eq:ip-fixed-a-forward}
\end{equation}
More precisely, for every \(M\geq1\), retain the leading and constant terms
and the Bernoulli terms with \(1\leq r\leq M-1\) (an empty sum when \(M=1\)).
The remainder is then \(\BigO(|t|^{2M-1})\), uniformly on the stated sets.
\end{theorem}

\begin{proof}
Insert, in \cref{eq:ip-lambert-log},
\[
 \frac1{1-\EulerE^{-x}}=\frac1x+\frac12+
 \sum_{r=1}^{M-1}\frac{B_{2r}}{(2r)!}x^{2r-1}+R_M(x).
\]
The displayed terms sum respectively to the dilogarithm, logarithm, and
polylogarithms in \cref{eq:ip-fixed-a-forward}.  To control the remainder,
choose a small \(\delta>0\) and split the Lambert index at
\(m\leq\delta/|t|\).  Taylor's theorem is uniform on this first range and
gives a contribution bounded by
\[
 C|t|^{2M-1}\sum_{m\geq1}\rho^m m^{2M-2}
 =\BigO(|t|^{2M-1}).
\]
On the complementary range, \(\rho^m\leq
\exp(-c/|t|)\).  Because \(t\) stays in a closed subsector of the right
half-plane, \(1-\EulerE^{-mt}\) is uniformly separated from zero there; the tail
is exponentially small and hence smaller than every power of \(|t|\).
\end{proof}

\begin{theorem}[Rigorous large-log base inverse]
\label{thm:ip-fixed-a-inverse}
Fix \(a\) with \(|a|<1\) and \(A=\Polylogarithm_2(a)\ne0\), and choose compatible
branches of all logarithms.  Put
\[
 D=\frac{\Polylogarithm_0(a)}{12},\qquad
 E=\frac{\Polylogarithm_{-2}(a)}{720},\qquad
 F=\frac{\Polylogarithm_{-4}(a)}{30240},\qquad
 X=-\log y+\frac12\log(1-a).
\]
Let \(X\to\infty\) so that \(A/X\) remains in a closed subsector of
\(\Re t>0\).  On the inverse sheet with \(t\sim A/X\),
\begin{align}
 t={}&\frac AX+\frac{DA^2}{X^3}
 +\frac{A^3(2D^2-AE)}{X^5}\notag\\
 &+\frac{A^4(A^2F-6ADE+5D^3)}{X^7}
 +\BigO(X^{-9}),\qquad q=\EulerE^{-t}.
 \label{eq:ip-fixed-a-inverse}
\end{align}
For real \(0<a<1\), this is the unique real base inverse as \(y\downarrow0\).
\end{theorem}

\begin{proof}
The preceding theorem gives
\[
 X=\frac At+Dt-Et^3+Ft^5+\BigO(t^7).
\]
Insert an ansatz \(t=A/X+\alpha X^{-3}+\beta X^{-5}+\gamma X^{-7}\).
Comparison of \(X^{-1},X^{-3},X^{-5}\) gives
\[
 \alpha=DA^2,\quad
 \beta=A^3(2D^2-AE),\quad
 \gamma=A^4(A^2F-6ADE+5D^3).
\]
This proves the formal coefficients, but the error also needs transfer.
For the displayed approximation \(t_0\), substitution in the forward
equation leaves a residual \(\BigO(X^{-7})\), while
\[
 \frac{\Differential}{\Differential t}\left(\frac At+Dt-Et^3+Ft^5\right)
 =-\frac A{t^2}+\BigO(1)
\]
has modulus comparable with \(|X|^2\).  The remainder is holomorphic and
uniformly \(\BigO(t^7)\) on a slightly wider closed subsector, so Cauchy's
estimate makes its derivative \(\BigO(t^6)\).  Taylor expansion of the exact
target equation on
\(|t-t_0|=K|X|^{-9}\) therefore consists of a constant
\(\BigO(X^{-7})\), a linear term of modulus comparable with
\(K|X|^{-7}\), and a quadratic error
\(\BigO(|X|^3|X|^{-18})\).  For fixed sufficiently large \(K\), the linear
term dominates the other two.  Rouché's theorem gives one zero inside that
circle, proving the \(\BigO(X^{-9})\) remainder.  The real uniqueness follows
from \cref{thm:ip-real-base}.
\end{proof}

\section{\texorpdfstring{The coalescing parameter \(a=q^x\)}{The coalescing parameter a equals q to the x}}
\label{sec:ip-qx}

\begin{theorem}[All-order coalescing expansion]
\label{thm:ip-qx-expansion}
Let \(\Re x>0\), \(q=\EulerE^{-t}\), and let \(t\to0\) in a closed subsector of
\(|\arg t|<\pi/2\).  Use the logarithm obtained by analytic continuation from
positive \(t\).  Then, locally uniformly for \(x\) in compact subsets of
\(\Re x>0\),
\begin{align}
 \log(q^x;q)_\infty
 \sim{}&-\frac{\pi^2}{6t}
 +\left(\frac12-x\right)\log t
 +\frac12\log(2\pi)-\log\Gamma(x) \notag\\
 &+\frac{B_2(x)}4t
 -\sum_{r\geq1}
 \frac{B_{2r}B_{2r+1}(x)}{2r(2r+1)!}t^{2r}.
 \label{eq:ip-qx-expansion}
\end{align}
\end{theorem}

\begin{proof}
For \(c>1\), initially with \(t>0\), Mellin inversion of the exponential in
the Lambert sum gives
\begin{equation}
 \log(\EulerE^{-xt};\EulerE^{-t})_\infty
 =-\frac1{2\pi\ImaginaryUnit}\int_{c-\ImaginaryUnit\infty}^{c+\ImaginaryUnit\infty}
 \Gamma(s)\zeta(s+1)\zeta(s,x)t^{-s}\Differential s.
 \label{eq:ip-qx-mellin}
\end{equation}
Absolute convergence for \(c>1\) justifies the interchange with the Lambert
sum.  Both sides then continue analytically to the stated sector, on the
logarithm sheet fixed in the theorem.  Indeed, uniformly for \(x\) in a fixed
compact subset of \(\Re x>0\), sufficiently small \(t\) in a smaller closed
subsector gives no factor zero: for bounded \(j\) the exponent
\((x+j)t\) has modulus below \(2\pi\), while for large \(j\) its argument
stays strictly between \(-\pi/2\) and \(\pi/2\), so it cannot be a nonzero
multiple of \(2\pi\ImaginaryUnit\).  We may therefore shift the continued contour left.
The pole at \(s=1\) contributes
\(-\pi^2/(6t)\).  At \(s=0\) the double pole, together with
\[
 \zeta(0,x)=\frac12-x,\qquad
 \zeta'(0,x)=\log\Gamma(x)-\frac12\log(2\pi),
\]
gives the logarithmic and constant terms.  At \(s=-1\), the residue of
\(\Gamma\), \(\zeta(0)=-1/2\), and
\(\zeta(-1,x)=-B_2(x)/2\) give \(B_2(x)t/4\).  At \(s=-2r\),
\[
 \operatorname*{Res}_{s=-2r}\Gamma(s)=\frac1{(2r)!},\quad
 \zeta(1-2r)=-\frac{B_{2r}}{2r},\quad
 \zeta(-2r,x)=-\frac{B_{2r+1}(x)}{2r+1},
\]
which yields the remaining coefficients, including the outer minus sign in
\cref{eq:ip-qx-mellin}.  At negative odd integers below \(-1\), a trivial zero
of \(\zeta(s+1)\) cancels the gamma pole.

Here is the required uniform remainder estimate.  Fix a compact set
\(K\Subset\{\Re x>0\}\), a closed subsector
\(|\arg t|\leq\pi/2-\delta\), and a vertical line \(\Re s=\sigma\) avoiding
the poles just listed.  Uniformly for \(x\in K\), Stirling's estimate and the
standard polynomial vertical bounds for the two zeta functions give constants
\(C,N\) such that
\[
 \left|\Gamma(\sigma+\ImaginaryUnit u)
 \zeta(\sigma+1+\ImaginaryUnit u)
 \zeta(\sigma+\ImaginaryUnit u,x)\right|
 \leq C(1+|u|)^N\EulerE^{-\pi|u|/2}.
\]
For the chosen branch of \(t^{-s}\),
\[
 |t^{-(\sigma+\ImaginaryUnit u)}|
 =|t|^{-\sigma}\EulerE^{u\arg t}
 \leq |t|^{-\sigma}\EulerE^{(\pi/2-\delta)|u|}.
\]
Thus the integrand on the shifted vertical line is bounded by
\[
 C|t|^{-\sigma}(1+|u|)^N\EulerE^{-\delta|u|},
\]
which is integrable uniformly for \(x\in K\).  The same exponential factor
makes the horizontal sides of the truncating rectangles tend to zero.  The
shifted contour is therefore \(\BigO_K(|t|^{-\sigma})\).  Shifting
successively past the displayed poles, and choosing \(\sigma\) strictly
between the last retained pole and the next uncancelled pole, proves the
asserted sectorial Poincaré remainder locally uniformly in \(x\).
\end{proof}

The specialization \(x=1\) has more structure than its power asymptotic:
Dedekind modularity restores the exponentially small remainder exactly, both
at the positive cusp and on the radial path toward \(-1\).

\begin{theorem}[Exact eta completions at \(q=1\) and \(q=-1\)]
\label{thm:ip-eta-completions}
For \(t>0\),
\begin{equation}
 (\EulerE^{-t};\EulerE^{-t})_\infty
 =\sqrt{\frac{2\pi}{t}}
 \exp\left(-\frac{\pi^2}{6t}+\frac{t}{24}\right)
 (\EulerE^{-4\pi^2/t};\EulerE^{-4\pi^2/t})_\infty.
 \label{eq:ip-eta-positive}
\end{equation}
If \(Q_c=\exp(-c\pi^2/t)\), then
\begin{equation}
 (-\EulerE^{-t};-\EulerE^{-t})_\infty
 =\sqrt{\frac{\pi}{t}}
 \exp\left(-\frac{\pi^2}{24t}+\frac{t}{24}\right)
 \frac{(Q_2;Q_2)_\infty^3}
 {(Q_4;Q_4)_\infty(Q_1;Q_1)_\infty}.
 \label{eq:ip-eta-negative}
\end{equation}
Consequently the relative errors after deleting the final product factor in
\cref{eq:ip-eta-positive} and the final product quotient in
\cref{eq:ip-eta-negative} are, respectively,
\(\BigO(\EulerE^{-4\pi^2/t})\) and \(\BigO(\EulerE^{-\pi^2/t})\).
\end{theorem}

\begin{proof}
Recall
\[
 \eta(\tau)=\EulerE^{\pi\ImaginaryUnit\tau/12}
 \prod_{n\geq1}(1-\EulerE^{2\pi\ImaginaryUnit n\tau})
\]
and the modular identity
\(\eta(-1/\tau)=\sqrt{-\ImaginaryUnit\tau}\,\eta(\tau)\), with the positive square root
when \(\tau\) is positive imaginary.  Set \(\tau=\ImaginaryUnit t/(2\pi)\).  Then
\begin{align*}
 \eta(\tau)&=\EulerE^{-t/24}(\EulerE^{-t};\EulerE^{-t})_\infty,\\
 \eta(-1/\tau)&=\EulerE^{-\pi^2/(6t)}
 (\EulerE^{-4\pi^2/t};\EulerE^{-4\pi^2/t})_\infty,\\
 \sqrt{-\ImaginaryUnit\tau}&=\sqrt{t/(2\pi)}.
\end{align*}
Solving the modular identity for the first product proves
\cref{eq:ip-eta-positive}.

For the negative radial path put \(x=\EulerE^{-t}\).  Separating the even and odd
powers in the defining product, and then separating the odd powers in the
ordinary positive product, gives
\begin{align*}
 (-x;-x)_\infty
  &=(x^2;x^2)_\infty(-x;x^2)_\infty,\\
 (-x;x^2)_\infty
  &=\frac{(x^2;x^4)_\infty}{(x;x^2)_\infty},\\
 (x;x^2)_\infty
  &=\frac{(x;x)_\infty}{(x^2;x^2)_\infty},
 \qquad
 (x^2;x^4)_\infty
  =\frac{(x^2;x^2)_\infty}{(x^4;x^4)_\infty}.
\end{align*}
Thus
\begin{equation}
 (-x;-x)_\infty
 =\frac{(x^2;x^2)_\infty^3}
 {(x;x)_\infty(x^4;x^4)_\infty}.
 \label{eq:ip-eta-dissection}
\end{equation}
Apply \cref{eq:ip-eta-positive} with \(t,2t,4t\).  The square-root factors
simplify to \(\sqrt{\pi/t}\), the elementary exponent to
\(-\pi^2/(24t)+t/24\), and the three complementary nomes to
\(Q_4,Q_2,Q_1\), respectively.  Substitution in
\cref{eq:ip-eta-dissection} proves \cref{eq:ip-eta-negative}.

Finally, as \(Q\downarrow0\), expansion of the first factor and the product
tail estimate \cref{eq:ip-product-tail} give
\((Q;Q)_\infty=1+\BigO(Q)\); the same estimate keeps its reciprocal bounded.
Applying this observation to the exact final factors proves both relative
error assertions.
\end{proof}

For real \(x>0\), the product decreases from \(1\) to \(0\) as \(q\) runs from
\(0\) to \(1\), since every factor \(1-q^{x+j}\) does.  It therefore has a
unique real base inverse.  Put
\[
 \Lambda=-\log y,\quad A=\frac{\pi^2}{6},\quad
 B=x-\frac12,\quad C=\log\Gamma(x)-\frac12\log(2\pi).
\]
The leading equation is \(\Lambda=A/t+B\log t+C\).  If \(B\ne0\), its exact
real Lambert-\(W\) solution for sufficiently large \(\Lambda\) is
\begin{equation}
 t_0=-\frac{A}
 {B\,W_k\!\left(-\frac AB
 \exp\left(-\frac{\Lambda-C}{B}\right)\right)},
 \label{eq:ip-qx-lambert}
\end{equation}
where \(k=-1\) when \(B>0\), and \(k=0\) when \(B<0\).
When \(B=0\), \(t_0=A/(\Lambda-C)\).  Exponentiating the leading equation and
putting \(w=-A/(Bt)\) proves \cref{eq:ip-qx-lambert}; all quantities are real,
so this exponentiation loses no logarithm-sheet information.  The omitted
part of \cref{eq:ip-qx-expansion} is \(\BigO(t)\), while the derivative of
\(A/t+B\log t+C\) has magnitude comparable with \(t^{-2}\).  The real
mean-value theorem, applied between the exact inverse and \(t_0\), therefore
gives \(t=t_0+\BigO(t_0^3)\).

\section{\texorpdfstring{Fixed-target double scaling at \(q=1\)}{Fixed-target double scaling at q=1}}
\label{sec:ip-double-scaling}

\begin{theorem}[All-order fixed-target inverse and recurrence]
\label{thm:ip-double-scaling}
Fix \(L>0\).  The unique real solution
\[
 (a;\EulerE^{-\varepsilon})_\infty=\EulerE^{-L},\qquad 0<a<1,
\]
has, for every \(M\), an asymptotic expansion
\(a(\varepsilon)=\sum_{k=1}^{M}A_k(L)\varepsilon^k+
\BigO(\varepsilon^{M+1})\).  Define
\[
 C_{r,n}=[\varepsilon^n]
 \left(\sum_{k\geq1}A_k\varepsilon^k\right)^r.
\]
Then \(A_1=L\), and for \(n\geq1\),
\begin{align}
 A_{n+1}={}&-\sum_{r=2}^{n+1}\frac{C_{r,n+1}}{r^2}
 -\frac12\sum_{r=1}^{n}\frac{C_{r,n}}r \notag\\
 &-\sum_{j=1}^{\Floor{n/2}}
 \frac{B_{2j}}{(2j)!}
 \sum_{r=1}^{n-2j+1}r^{2j-2}C_{r,n-2j+1}.
 \label{eq:ip-double-scaling-recurrence}
\end{align}
In particular,
\begin{align}
 a={}&L\varepsilon-\frac{L(L+2)}4\varepsilon^2
 +\frac{L(L^2+9L+12)}{72}\varepsilon^3 \notag\\
 &-\frac{L(L^3+4L^2+12L+24)}{576}\varepsilon^4 \notag\\
 &-\frac{L(31L^4-75L^3-800L^2-720)}
 {86400}\varepsilon^5+\BigO(\varepsilon^6).
 \label{eq:ip-double-scaling-five}
\end{align}
Every \(A_k\) is a rational polynomial in \(L\) divisible by \(L\).
\end{theorem}

\begin{proof}
Taking the negative logarithm gives the exact positive Lambert equation
\begin{equation}
 \Phi_\varepsilon(a):=
 \sum_{r\geq1}\frac{a^r}{r(1-\EulerE^{-r\varepsilon})}=L.
 \label{eq:ip-positive-lambert}
\end{equation}
Since the first summand is at least \(a/\varepsilon\), every solution obeys
\(0<a\leq L\varepsilon\).  Moreover
\[
 \partial_a\Phi_\varepsilon(a)
 =\sum_{r\geq1}\frac{a^{r-1}}{1-\EulerE^{-r\varepsilon}}
 \geq\frac1{1-\EulerE^{-\varepsilon}}\geq\frac1\varepsilon.
\]
Thus \(\Phi_\varepsilon\) is strictly increasing, and existence and uniqueness
also follow from the principal real argument atlas.

Expand
\[
 \frac1{1-\EulerE^{-r\varepsilon}}
 =\frac1{r\varepsilon}+\frac12+
 \sum_{j\geq1}\frac{B_{2j}}{(2j)!}
 (r\varepsilon)^{2j-1}.
\]
The coefficient of \(\varepsilon^n\) in
\cref{eq:ip-positive-lambert} contains the new unknown \(A_{n+1}\) only in
\(a/\varepsilon\), with coefficient one.  The remaining contributions are,
respectively, the three sums on the right of
\cref{eq:ip-double-scaling-recurrence}; setting the coefficient to zero proves
the recurrence.  It gives the five displayed polynomials by successive
substitution.  Inductively its operations preserve rational polynomiality,
and every term vanishes at \(L=0\), proving the factor \(L\).

It remains to justify the asymptotic remainder.  For a truncation
\(a_M=\sum_{k=1}^{M}A_k\varepsilon^k=\BigO(\varepsilon)\), use
\[
 1-\EulerE^{-r\varepsilon}\geq c\min(r\varepsilon,1)
\]
for real \(\varepsilon>0\).  Splitting at \(r=\varepsilon^{-1}\) then bounds
the Lambert tail \(r\geq M+1\) by
\[
 \frac C\varepsilon
 \sum_{M+1\leq r\leq\varepsilon^{-1}}\frac{|a_M|^r}{r^2}
 +C\sum_{r>\varepsilon^{-1}}\frac{|a_M|^r}{r}
 =\BigO(\varepsilon^M).
\]
For the finitely many retained indices, Taylor's theorem for the elementary
kernel, together with the recurrence, gives
\(\Phi_\varepsilon(a_M)-L=\BigO(\varepsilon^M)\).  The derivative lower bound
and the real mean-value theorem then imply
\[
 |a-a_M|\leq\varepsilon\,|\Phi_\varepsilon(a)-\Phi_\varepsilon(a_M)|
 =\BigO(\varepsilon^{M+1}).
\]
This completes both the formal and analytic parts of the proof.
\end{proof}

\begin{warning}[Complex double scaling]
For complex \(L=-\log y\), the recurrence still defines a formal branch, but
the poles \(2\pi\ImaginaryUnit k/r\) of the Lambert kernels accumulate at
\(\varepsilon=0\).  The real theorem above does not imply a convergent complex
Taylor series.  A complex version requires a sector, a logarithm branch, and
uniform remainder estimates; absent those data it remains a formal
asymptotic construction.
\end{warning}

\section{Raw root-of-unity endpoints and radial inverses}
\label{sec:ip-cyclotomic}

\begin{proposition}[Raw boundary partial products]
\label{prop:ip-boundary-products}
Let \(\zeta\) be primitive of order \(d\), and write \(N=kd+s\),
\(0\leq s<d\).  Then
\begin{equation}
 (a;\zeta)_N=(1-a^d)^k\prod_{j=0}^{s-1}(1-a\zeta^j).
 \label{eq:ip-boundary-block}
\end{equation}
Hence the partial products tend to zero when \(|1-a^d|<1\), are unbounded
when \(|1-a^d|>1\), and do not converge when \(|1-a^d|=1\) and \(a\ne0\).
For \(a=0\), they are identically one.
\end{proposition}

\begin{proof}
Every complete residue block contributes
\(\prod_{j=0}^{d-1}(1-a\zeta^j)=1-a^d\), proving
\cref{eq:ip-boundary-block}.  The first two conclusions follow from the
geometric factor.  In the unit-modulus case, the subsequence \(s=0\) is
\((1-a^d)^k\); it could converge only if \(1-a^d=1\), which forces \(a=0\).
\end{proof}

This proposition concerns raw partial products on the boundary.  The next
theorem instead approaches the boundary from \(|q|<1\), where the analytic
product is defined.

\begin{theorem}[Uniform radial cyclotomic expansion and inverse]
\label{thm:ip-radial-cyclotomic}
Let \(\zeta\) be primitive of order \(d\), let \(|a|\leq\rho<1\), and put
\(q=\zeta\EulerE^{-t}\), with \(t\to0\) in a closed subsector of \(\Re t>0\).
All logarithms below are the holomorphic determinations obtained by
continuation from \(a=0\); in particular,
\[
 \log(1-a^d)=-\sum_{\ell\geq1}\frac{a^{d\ell}}{\ell}
 \qquad (|a|<1).
\]
Define
\begin{align}
 A_\zeta(a)&=\frac{\Polylogarithm_2(a^d)}{d^2},
 \label{eq:ip-cyclotomic-A}\\
 C_\zeta(a)&=\frac1{2d}\log(1-a^d)
 -\sum_{\substack{m\geq1\\d\nmid m}}
 \frac{a^m}{m(1-\zeta^m)},
 \label{eq:ip-cyclotomic-C}\\
 D_\zeta(a)&=\frac{a^d}{12(1-a^d)}
 -\sum_{\substack{m\geq1\\d\nmid m}}
 \frac{a^m\zeta^m}{(1-\zeta^m)^2}.
 \label{eq:ip-cyclotomic-D}
\end{align}
The coefficient series converge absolutely, and
\begin{equation}
 \log(a;\zeta\EulerE^{-t})_\infty
 =-\frac{A_\zeta(a)}t+C_\zeta(a)-D_\zeta(a)t
 +\BigO(t^2)
 \label{eq:ip-cyclotomic-forward}
\end{equation}
uniformly on the stated compact sets.  If \(A_\zeta(a)\ne0\), choose a branch
of \(\log y\), put \(X=-\log y+C_\zeta(a)\), and suppose that \(X\to\infty\)
with \(A_\zeta(a)/X\) in the chosen closed subsector.  The corresponding
inverse satisfies
\begin{equation}
 t=\frac{A_\zeta(a)}X+
 \frac{D_\zeta(a)A_\zeta(a)^2}{X^3}
 +\BigO(X^{-4}),\qquad q=\zeta\EulerE^{-t}.
 \label{eq:ip-cyclotomic-inverse}
\end{equation}
\end{theorem}

\begin{proof}
Split \cref{eq:ip-lambert-log} into indices \(m=d\ell\) and \(d\nmid m\).
For the multiples,
\[
 \frac1{1-\EulerE^{-d\ell t}}
 =\frac1{d\ell t}+\frac12+\frac{d\ell t}{12}
 +\BigO((\ell t)^3).
\]
Their singular, constant, and linear sums are respectively
\[
 -\frac{\Polylogarithm_2(a^d)}{d^2t},\qquad
 \frac1{2d}\log(1-a^d),\qquad
 -\frac{a^d}{12(1-a^d)}t.
\]
For nonmultiples,
\[
 \frac1{1-\zeta^m\EulerE^{-mt}}
 =\frac1{1-\zeta^m}
 -\frac{m\zeta^m}{(1-\zeta^m)^2}t
 +\BigO(m^2t^2).
\]
These are exactly the remaining terms in
\cref{eq:ip-cyclotomic-C,eq:ip-cyclotomic-D}.

Absolute convergence of the coefficient sums is not by itself a remainder
proof, so split again at \(m\leq\delta/|t|\).  Choose \(\delta\) below the
nearest singularity of each of the finitely many residue-class kernels.  On
the first range Taylor's theorem gives a uniform remainder, and summation of
\(\rho^m m|t|^2\) gives \(\BigO(t^2)\).  On the tail,
\(\rho^m=\BigO(\exp(-c/|t|))\); the sector condition also separates every
denominator from zero.  This proves \cref{eq:ip-cyclotomic-forward}.

The target equation is
\[
 X=\frac{A_\zeta}t+D_\zeta t+\BigO(t^2).
\]
Substitution of \(t_0=A_\zeta/X+D_\zeta A_\zeta^2/X^3\) leaves residual
\(\BigO(X^{-2})\), while the derivative has size comparable with \(X^2\).
Uniform holomorphy on a slightly wider subsector and Cauchy's estimate control
the derivatives of the \(\BigO(t^2)\) remainder.  Thus, on a circle of radius
\(K|X|^{-4}\) around \(t_0\), the linear term has size comparable with
\(K|X|^{-2}\), the constant residual is \(\BigO(X^{-2})\), and all higher
variations are smaller.  For fixed sufficiently large \(K\), Rouché's theorem
supplies a unique nearby solution and proves
\cref{eq:ip-cyclotomic-inverse}.
\end{proof}

\begin{corollary}[The radial points \(1\) and \(-1\)]
\label{cor:ip-plus-minus-one}
For real \(0<a<1\), \cref{eq:ip-fixed-a-inverse} is the unique inverse toward
\(q=1^-\).  For real \(0<|a|<1\), put
\[
 A_-=\frac14\Polylogarithm_2(a^2),\quad
 C_-=\frac12\log(1-a),\quad
 D_-=\frac{a(a+3)}{12(1-a^2)},
\]
\[
 E_-=\frac{a(15a^4+4a^3+90a^2+4a+15)}
 {720(1-a^2)^3},\qquad X=-\log y+C_-.
\]
On the real radial branch \(q=-\EulerE^{-t}\), \(t\downarrow0\),
\begin{align}
 \log(a;-\EulerE^{-t})_\infty
 &=-\frac{A_-}t+C_--D_-t+E_-t^3+\BigO(t^5),
 \label{eq:ip-minus-one-forward}\\
 t&=\frac{A_-}X+\frac{D_-A_-^2}{X^3}
 +\frac{2D_-^2A_-^3-E_-A_-^4}{X^5}
 +\BigO(X^{-7}).
 \label{eq:ip-minus-one-inverse}
\end{align}
Here the inverse formula is taken as
\(y=(a;-\EulerE^{-t})_\infty\downarrow0\), equivalently \(X\to+\infty\).
\end{corollary}

\begin{proof}
Split the canonical Lambert logarithm into even and odd indices:
\[
 \log(a;-\EulerE^{-t})_\infty
 =-\sum_{m\geq1}\frac{a^m}{m(1-(-1)^m\EulerE^{-mt})}.
\]
For even indices use
\[
 \frac1{1-\EulerE^{-x}}
 =\frac1x+\frac12+\frac{x}{12}-\frac{x^3}{720}
 +\BigO(x^5),
\]
and for odd indices use
\[
 \frac1{1+\EulerE^{-x}}
 =\frac12+\frac{x}{4}-\frac{x^3}{48}+\BigO(x^5).
\]
Only the even indices \(m=2r\) contribute a singular term, namely
\[
 -\sum_{r\geq1}\frac{a^{2r}}{2r}\frac1{2rt}
 =-\frac{\Polylogarithm_2(a^2)}{4t}.
\]
The constant term is
\[
 -\frac12\sum_{m\geq1}\frac{a^m}{m}
 =\frac12\log(1-a).
\]
The coefficient of \(t\) is
\[
 -\frac1{12}\sum_{r\geq1}a^{2r}
 -\frac14\sum_{\substack{m\geq1\\m\ \mathrm{odd}}}a^m
 =-\frac{a(a+3)}{12(1-a^2)}.
\]
Writing \(z=a^2\), the coefficient of \(t^3\) is
\[
 \frac1{180}\sum_{r\geq1}r^2z^r
 +\frac1{48}\sum_{\substack{m\geq1\\m\ \mathrm{odd}}}m^2a^m.
\]
The elementary identities
\[
 \sum_{r\geq1}r^2z^r=\frac{z(1+z)}{(1-z)^3},
 \qquad
 \sum_{\substack{m\geq1\\m\ \mathrm{odd}}}m^2a^m
 =\frac{a(1+a)}{(1-a)^3}
 -\frac{4a^2(1+a^2)}{(1-a^2)^3}
\]
simplify this coefficient to
\[
 \frac{a(15a^4+4a^3+90a^2+4a+15)}{720(1-a^2)^3}=E_-.
\]
For fixed \(|a|<1\), the same finite-index/tail split used in the proof of
\cref{thm:ip-radial-cyclotomic} justifies the expansions under the sum; the
summed fifth-order remainders are \(\BigO(t^5)\).  This proves
\cref{eq:ip-minus-one-forward} with every displayed coefficient.

Now
\[
 X=\frac{A_-}{t}+D_-t-E_-t^3+\BigO(t^5).
\]
Substituting
\(t=A_-/X+\alpha X^{-3}+\beta X^{-5}\) and comparing the coefficients of
\(X^{-1}\) and \(X^{-3}\) gives
\[
 \alpha=D_-A_-^2,
 \qquad
 \beta=2D_-^2A_-^3-E_-A_-^4.
\]
The resulting residual is \(\BigO(X^{-5})\), whereas the derivative of the
target equation is
\(-A_-/t^2+\BigO(1)\asymp-X^2\).  The real mean-value theorem therefore
improves the inverse error to \(\BigO(X^{-7})\), proving
\cref{eq:ip-minus-one-inverse}.  The exclusion \(a=0\) is necessary: there
the product is identically one and \(A_-=0\).
\end{proof}

\begin{proposition}[Fixed-level cyclotomic coalescence with winding]
\label{prop:ip-cyclotomic-fixed-level}
Let \(\zeta\) be primitive of order \(d\), let
\(q=\zeta\EulerE^{-t}\), and let \(t\to0\) in a closed subsector of
\(\Re t>0\).  Fix \(y\in\ComplexNumbers^\times\), choose a logarithm
\(\ell\) of \(y\), and, for \(\nu\in\IntegerNumbers\), put
\[
 L_\nu:=-\ell-2\pi\ImaginaryUnit\nu.
\]
Suppose \(L_\nu\ne0\).  For the Lambert logarithm that is zero at \(a=0\),
the lifted equation
\[
 \log(a;q)_\infty=\ell+2\pi\ImaginaryUnit\nu
\]
has exactly \(d\) solution germs tending to \(a=0\).  They satisfy
\begin{equation}
 a^d\sim d^2L_\nu t,\qquad
 a\sim\omega(d^2L_\nu t)^{1/d},\qquad \omega^d=1.
 \label{eq:ip-cyclotomic-fixed-level}
\end{equation}
Consequently the raw equation \((a;q)_\infty=y\) has one such
\(d\)-sheeted cluster for every fixed winding \(\nu\) with \(L_\nu\ne0\).
Every continuous raw solution branch tending to zero belongs to exactly one
of these clusters.  The exceptional lift \(L_\nu=0\), when it occurs, is a
different degenerate problem and is not covered by the displayed scale.
\end{proposition}

\begin{proof}
Write \(\Lambda(a,t)\) for the Lambert logarithm in
\cref{eq:ip-cyclotomic-forward}.  The lifted equation is
\(-\Lambda(a,t)=L_\nu\).  Suppose first that a solution germ has
\(a\to0\).  In the uniform forward formula,
\(C_\zeta(a)=\LittleO(1)\) and \(D_\zeta(a)t=\LittleO(1)\); its remainder is
uniform for \(|a|\leq\rho\), so
\[
 \frac{\Polylogarithm_2(a^d)}{d^2t}=L_\nu+\LittleO(1).
\]
Since \(\Polylogarithm_2(z)/z\to1\), the last identity first gives
\(a^d=\BigO(t)\) and then
\(a^d/(d^2t)\to L_\nu\), proving the first relation and the scale
\(a=\BigO(t^{1/d})\).

Choose a determination \(\tau=t^{1/d}\) on the subsector and write
\(a=\tau b\).  Uniformly for \(b\) in a fixed compact set,
\[
 -\Lambda(\tau b,t)-L_\nu
 =\frac{b^d}{d^2}-L_\nu+\LittleO(1),
\]
because
\(\Polylogarithm_2(tb^d)/(d^2t)=b^d/d^2+\BigO(t)\), while
\(C_\zeta(\tau b)=\LittleO(1)\) and
\(tD_\zeta(\tau b)=\LittleO(1)\).  The limiting polynomial has the
\(d\) distinct roots of \(b^d=d^2L_\nu\).  Rouché's theorem on disjoint
small circles about those roots supplies one zero on each circle for all
sufficiently small \(t\).  The preceding converse estimate shows that every
solution tending to zero lies on one of these circles.  Each circle contains
one zero counted with multiplicity, so that zero is simple.  For every
nonzero \(t\) in the subsector, the holomorphic implicit-function theorem
therefore makes it a locally holomorphic function of the chosen coordinate
\(\tau=t^{1/d}\).  The circles are fixed and pairwise disjoint, and each
contains exactly one zero; uniqueness makes these local functions agree on
overlaps and patch to \(d\) sectorial solution germs.  Conversely, any
solution with \(a\to0\) has
\[
 b^d=d^2L_\nu+\LittleO(1),\qquad b=a/\tau,
\]
and hence eventually lies in exactly one of these circles.  This proves the
exact germ count, its exhaustivity, and both asymptotic relations.

Finally, exponentiating a lifted solution gives the raw equation.  Conversely,
on a continuous raw solution branch the quantity
\((\Lambda(a,t)-\ell)/(2\pi\ImaginaryUnit)\) is integer-valued and continuous,
hence is one fixed integer \(\nu\).  Thus the winding clusters are disjoint
and exhaust all continuous raw branches tending to zero.
\end{proof}

\begin{warning}[Exterior normalization]
The defining product does not converge for \(|q|>1\).  There is therefore no
canonical unnormalized value of \((a;q)_\infty\), and no canonical inverse at
\(q=\infty\).  One may instead specify, for example,
\((a;q^{-1})_\infty\) for \(|q|>1\), or a \(q\)-gamma normalization with its
explicit powers removed.  Different choices can have different inverse
asymptotics and must not be conflated.
\end{warning}
